\documentclass[8pt, aspectratio=169]{beamer}
%%%%%%%%%%%%%%%%%%%%%%%%%%%%%%%%%
% PACKAGE IMPORTS
%%%%%%%%%%%%%%%%%%%%%%%%%%%%%%%%%


%\usepackage[tmargin=2cm,bmargin=2cm,footskip=.2in, rmargin=1.5in,lmargin=1.5in, margin=0.85in, a4paper]{geometry}
\usepackage{bookmark}
% \usepackage{libertinus}
\usepackage{amsmath,amsfonts,amsthm,amssymb,mathtools}
%\usepackage[varbb]{newpxmath}
\usepackage[libertine]{newtxmath}
\usepackage{chemfig}
\usepackage{ragged2e}
\usepackage{xfrac}
\usepackage{mhchem}
\usepackage[italian]{babel}
% \usepackage[Sonny]{fncychap}
\usepackage[makeroom]{cancel}
\usepackage{mathtools}
\usepackage{listing}
\usepackage{enumitem}
\usepackage{hyperref,theoremref}
\hypersetup{
	pdftitle={Espansione in multipolo del campo gravitazionale newtoniano},
	colorlinks=true, linkcolor=doc!90,
	bookmarksnumbered=true,
	bookmarksopen=true,
	pdfauthor={Francesco Sermi}
}
\usepackage[most,many,breakable]{tcolorbox}
\usepackage{xcolor}
\usepackage{varwidth}
\usepackage{varwidth}
\usepackage{etoolbox}
%\usepackage{authblk}
\usepackage{nameref}
\usepackage{multicol,array}
\usepackage{tikz}
\usepackage{tikz-cd}
\usepackage{tikz-3dplot}
\usepackage{tikz}
\usetikzlibrary{3d,decorations.pathmorphing}
\usepackage[ruled,vlined,linesnumbered]{algorithm2e}
\usepackage{import}
\usepackage{xifthen}
\usepackage{pdfpages}
\usepackage{transparent}

\newcommand\mycommfont[1]{\footnotesize\ttfamily\textcolor{blue}{#1}}
\SetCommentSty{mycommfont}
\newcommand{\incfig}[1]{%
    \def\svgwidth{\columnwidth}
    \import{./figures/}{#1.pdf_tex}
}

\usepackage{tikzsymbols}
\usepackage{float}
\usepackage[toc]{appendix}
\renewcommand\qedsymbol{$\square$}
\usepackage{hyperref}
\usepackage{mwe}
\usepackage{fancyhdr}
\usepackage{pgfplots}
\usepackage{bm}
\usepackage{dsfont}
%\usepackage{import}
%\usepackage{xifthen}
%\usepackage{pdfpages}
%\usepackage{transparent}
\definecolor{doc}{RGB}{0,0,0}

\newcommand*\circled[1]{\tikz[baseline=(char.base)]{
		\node[shape=circle,draw,inner sep=1pt] (char) {#1};}}
\newcommand{\innerprod}[2]{\langle #1, #2 \rangle}
\renewcommand\appendixpagename{Appendici}
\renewcommand\appendixtocname{Appendici}
\setcounter{localmathalphabets}{0}

\usepackage{amsmath,amssymb}

\setbeamertemplate{frame numbering}{}
\definecolor{accent}{HTML}{2E5EAA}
\setbeamercolor{palette primary}{bg=accent,fg=white}
\setbeamercolor{frametitle}{bg=accent,fg=white}
\setbeamercolor{progress bar}{fg=accent,bg=accent!20}

\title{Espansione in multipoli della gravità}
\author{Francesco Sermi}
\institute{Relatore: Prof. Leonardo Gualtieri}
\date{Anno Accademico 2025/2026}

\begin{document}

\maketitle

\begin{frame}{Equazioni della gravità newtoniana}
	Le equazioni del moto situato in $\vec{r}(t)$ di un corpo massivo con massa inerziale $m_I$ e massa gravitazionale $m_G$ soggetto alla forza gravitazionale generata da un altro corpo con massa $M$ posto nell'origine del nostro sistema di riferimento sono date da
$$
m_I \, \vec{a}(t) = -G\, m_G M \, \frac{\vec{r}(t)}{|\vec{r}(t)|^3},
$$
dove $G$ è la costante di gravitazione universale. Introducendo il \emph{potenziale gravitazionale}
$$
	U(\vec{r}) = -\frac{GM}{|\vec{r}(t)|} \implies m_I \vec{a}(t) = -m_G \vec{\nabla} U(\vec{r})
$$
e supponendo che
\begin{enumerate}
	\item valga il \textbf{principio di equivalenza debole}, ovvero $m_I = m_G$;
	\item valga il \textbf{principio di sovrapposizione}, ossia il potenziale in un punto è dato dalla sovrapposizione lineare dei potenziali gravitazionali generati dagli altri corpi presenti nel sistema studiato
	$$
		U(\vec{x}) = \sum_i^N U_i(\vec{x})
	$$
\end{enumerate}
\end{frame}
\begin{frame}{Sviluppo in multipoli}
Usando il principio di sovrapposizione, è facile dimostrare che per una distribuzione continua di massa caratterizzato da una certa distribuzione volumica di massa $\rho(\vec{x},t)$, allora il potenziale in un punto è dato da
\begin{align*}
	U = -G\int_V \frac{\rho(\vec{x}', t)}{|\vec{x} - \vec{x}'|} \mathrm{d}^3 x \mathrm{d}^3 x'
\end{align*}
e usando la ben nota identità $\nabla^2 \frac{1}{|\vec{x} - \vec{x}'|} = -4 \pi \delta^3(\vec{x} - \vec{x}')$ si ricava che il potenziale gravitazionale soddisfa l'equazione di Poisson
\begin{align*}
	\nabla^2 U = 4 \pi G \rho(\vec{x}, t)
\end{align*}
Per corpi 
\end{frame}

\begin{frame}{Definizione dei Momenti di Massa}
I momenti di massa e di corrente racchiudono tutte le proprietà fisiche della sorgente estesa.

\bigskip
Nel formalismo STF, il momento di massa di ordine $l$ è definito come l'integrale della densità di massa pesata sui prodotti tensoriali delle coordinate spaziali.
\end{frame}

\begin{frame}{La Deformabilità Mareale}
La deformabilità mareale quantifica la risposta strutturale di un corpo esteso all'influenza di un campo gravitazionale esterno.

\bigskip
Nel limite di campo debole, il momento di quadripolo indotto è linearmente proporzionale al campo di marea esterno applicato, tramite una costante definita deformabilità mareale (strettamente legata al numero di Love di ordine 2).
\end{frame}

\begin{frame}{Approccio Classico vs. Modellizzazione Relativistica}
\begin{columns}[T,onlytextwidth]
\column{0.48\textwidth}
\textbf{Gravità Newtoniana}\\[0.5em]
In ambito classico, l'analisi si semplifica notevolmente. Il problema fisico si riduce alla risoluzione dell'equazione di Poisson per il potenziale perturbato.

\column{0.48\textwidth}
\textbf{Relatività Generale}\\[0.5em]
La determinazione della deformabilità richiede la complessa integrazione delle equazioni di Einstein perturbate, accoppiando in modo coerente la metrica dello spaziotempo all'equazione di stato della materia interna.
\end{columns}
\end{frame}

\end{document}
